\chapter*{中期进展概述}
\addcontentsline{toc}{chapter}{中期进展概述}

本课题面向海上风电电缆运维，研究将多源感知、状态估计、任务决策与存量AUV执行平台衔接的安全自主巡检方法。本报告以\ReportCutoff 为事实截止日，按研究问题、必要方法、代表结果和后续验证路径组织阶段工作。

感知方面，已建立先验区域搜索、磁检测、主动横切、显式半高全宽估计和沿缆跟踪流程，并以稳定性对照识别增益变化、场形失配与目标关联的不同影响。条件性运动仿真检验了横向巡检及拒绝机制，深度发布则与场形和高度依据分开管理。估计和控制方面，已完成异常观测处理与不确定性调制对照，明确平滑性、精度和约束满足之间的权衡。工程方面，磁测量链、Jetson独立求解器和真实PC104通信与安全状态链为方法提供实物子链支撑。

现阶段工作的主要转向是将已形成的算法与接口证据连接到同一实物平台。能源、必要反馈、执行、装配水密及水中受控运行仍需逐级确认；磁载荷装艇、目标平台完整时序和真实海缆任务随后推进。报告将这些工作按依赖和验收证据安排，不以子链通过替代整艇完成。

\begin{table}[htbp]
\centering\caption{报告阅读入口与阶段状态}
\begin{tabularx}{\linewidth}{p{2.5cm}Xp{2.2cm}}\toprule
研究问题&已有进展与下一验证&正文位置\\\midrule
磁几何与目标&剖面对照及条件巡检已验证；真实场形/身份/高度依据待完善&第3章\\
融合与运动控制&已获得对照及性能权衡；真实传感与执行约束待验证&第4章\\
测量与运行接口&已形成实物子链；装艇后耦合与完整时序待验证&第4章\\
整艇与完整任务&基础恢复、水中受控运行及任务级实物验证待完成&第5、6章\\\bottomrule
\end{tabularx}
\end{table}
